% Part: sets-functions-relations
% Chapter: functions
% Section: kinds

\documentclass[../../../include/open-logic-section]{subfiles}

\begin{document}

\olfileid{sfr}{fun}{kin}
\olsection{অপেক্ষকের প্রকারভেদ}

\begin{explain}
যেসব ধরনের অপেক্ষক আমাদের সবচেয়ে বেশি কাজে আসে, সেগুলির একটি শ্রেণিবিন্যাস প্রবর্তন করা উপযোগী হবে।

প্রথমে এমন অপেক্ষক বিবেচনা করা যাক, যার সহসংজ্ঞাক্ষেত্রের প্রত্যেক সদস্যই অপেক্ষকটির একটি মান। এই ধরনের অপেক্ষককে !!{surjective} বলা হয়; \olref{fig:surjective}-এর মতো ছবি দিয়ে তাদের দেখানো যায়।

\begin{figure}
  \olasset{assets/diagrams/surjective.tikz}
  \caption{!!^a{surjective} অপেক্ষকের সহসংজ্ঞাক্ষেত্রের প্রত্যেক !!{element} অপেক্ষকটির একটি মান।}
  \ollabel{fig:surjective}
\end{figure}
\end{explain}

\begin{defn}[!!^{surjective} অপেক্ষক]
একটি অপেক্ষক $f \colon A \rightarrow B$ \emph{!!{surjective}} হয় যদি এবং কেবল যদি $B$ একই সঙ্গে~$f$-এর বিস্তৃতিও হয়, অর্থাৎ প্রত্যেক $y \in B$-এর জন্য অন্তত একটি $x \in A$ থাকে, যাতে~$f(x) = y$; প্রতীকে:
\[
  (\forall y \in B)(\exists x \in A)f(x) = y.
\]
এই ধরনের অপেক্ষককে $A$ থেকে $B$-তে !!a{surjection} বলা হয়।
\end{defn}

\begin{explain}
$f$ !!a{surjection} প্রমাণ করতে হলে দেখাতে হবে যে $f$-এর সহসংজ্ঞাক্ষেত্রের প্রত্যেক বস্তু কোনো ইনপুট $x$-এর জন্য $f(x)$-এর মান।

লক্ষ করো, যেকোনো অপেক্ষক !!a{surjection} \emph{নির্ধারণ করে}। বস্তুত, একটি অপেক্ষক $f \colon A \to B$ দেওয়া থাকলে, $f'(x) = f(x)$ দ্বারা $f' \colon A \to \ran{f}$ সংজ্ঞায়িত করি। যেহেতু $\ran{f}$-কে $\Setabs{f(x) \in B}{x \in A}$ হিসেবে \emph{সংজ্ঞায়িত} করা হয়, তাই এই অপেক্ষক $f'$ অবশ্যই !!a{surjection}।
\end{explain}

\begin{explain}
যেকোনো অপেক্ষক প্রতিটি সম্ভাব্য ইনপুটকে একটিমাত্র আউটপুটে পাঠায়। কিন্তু এমন অপেক্ষকও আছে, যারা কখনো ভিন্ন ইনপুটকে একই আউটপুটে পাঠায় না। এই ধরনের অপেক্ষককে !!{injective} বলা হয়; \olref{fig:injective}-এর মতো ছবি দিয়ে তাদের দেখানো যায়।
\begin{figure}
  \olasset{assets/diagrams/injective.tikz}
  \caption{!!^a{injective} অপেক্ষক কখনো দুটি ভিন্ন আর্গুমেন্টকে একই মানে পাঠায় না।}
  \ollabel{fig:injective}
\end{figure}
\end{explain}

\begin{defn}[!!^{injective} অপেক্ষক]
একটি অপেক্ষক $f \colon A \rightarrow B$ \emph{!!{injective}} হয় যদি এবং কেবল যদি প্রত্যেক $y \in B$-এর জন্য বড়জোর একটি $x \in A$ থাকে, যাতে~$f(x) = y$। এই ধরনের অপেক্ষককে $A$ থেকে~$B$-তে !!a{injection} বলা হয়।
\end{defn}

\begin{explain}
$f$ !!a{injection} প্রমাণ করতে হলে দেখাতে হবে যে $f$-এর সংজ্ঞাক্ষেত্রের যেকোনো !!{element}s $x$ ও $y$-এর ক্ষেত্রে, $f(x)=f(y)$ হলে $x=y$।
\end{explain}

\begin{ex}
$f(x) = 1$ দ্বারা নির্ধারিত ধ্রুব অপেক্ষক $f\colon \Nat \to \Nat$ না !!{injective}, না !!{surjective}।

$f(x) = x$ দ্বারা নির্ধারিত অভেদ অপেক্ষক $f\colon \Nat \to \Nat$ একই সঙ্গে !!{injective} ও !!{surjective}।

$f(x) = x+1$ দ্বারা নির্ধারিত উত্তরসূরি অপেক্ষক $f \colon \Nat \to \Nat$ !!{injective}, কিন্তু !!{surjective} নয়।

নিচের সংজ্ঞা দ্বারা নির্ধারিত অপেক্ষক $f \colon \Nat \to \Nat$:
\[
  f(x) =
  \begin{cases}
    \frac{x}{2} & \text{যদি $x$ জোড় হয়} \\
    \frac{x+1}{2} & \text{যদি $x$ বিজোড় হয়।}
  \end{cases}
\]
!!{surjective}, কিন্তু !!{injective} নয়।
\end{ex}

\begin{explain}
প্রায়ই আমরা এমন অপেক্ষক বিবেচনা করতে চাই, যা একই সঙ্গে !!{injective} ও !!{surjective}। এই ধরনের অপেক্ষককে !!{bijective} বলা হয়। তাদের চেহারা \olref{fig:bijective}-এ দেখানো অপেক্ষকটির মতো। !!^{bijection}s-কে কখনো কখনো \emph{এক-এক সঙ্গতি}ও বলা হয়, কারণ তারা সহসংজ্ঞাক্ষেত্রের উপাদানগুলিকে সংজ্ঞাক্ষেত্রের উপাদানগুলির সঙ্গে অনন্যভাবে জোড়া মেলায়।
\begin{figure}
  \olasset{assets/diagrams/bijective.tikz}
  \caption{!!^a{bijective} অপেক্ষক সহসংজ্ঞাক্ষেত্রের উপাদানগুলিকে সংজ্ঞাক্ষেত্রের উপাদানগুলির সঙ্গে অনন্যভাবে জোড়া মেলায়।}
  \ollabel{fig:bijective}
\end{figure}
\end{explain}

\begin{defn}[!!^{bijection}]
একটি অপেক্ষক $f \colon A \to B$ \emph{!!{bijective}} হয় যদি এবং কেবল যদি সেটি একই সঙ্গে !!{surjective} ও !!{injective} হয়। এই ধরনের অপেক্ষককে $A$ থেকে~$B$-তে (বা $A$ ও~$B$-এর মধ্যে) !!a{bijection} বলা হয়।
\end{defn}

\end{document}
